% Source: https://artofproblemsolving.com/community/c3h1410836_pigeonhole_principle
% Stable identity: aops-topic-1410836-q01
\textbf{Answer: $505$.}

\textbf{Construction.} Take $M=\{1,5,9,\dots,2017\}$, the numbers congruent to $1$ modulo $4$. Any two differ by a multiple of $4$, which is not prime. There are $\dfrac{2017-1}{4}+1=505$ of them.

\textbf{Bound.} Among any $8$ consecutive integers $a,a+1,\dots,a+7$, the set $M$ has at most $2$ elements. Indeed, two elements of such a block differ by $1,2,\dots,7$, and the non-prime values among these are $1,4,6$. If $x<y<z$ were three elements of $M$ in the block, then $y-x$, $z-y$ and $z-x=(y-x)+(z-y)$ would all lie in $\{1,4,6\}$. But no two elements of $\{1,4,6\}$ (possibly equal) add up to an element of $\{1,4,6\}$: the possible sums are $2,5,7,8,10,12$. Contradiction.

Now split $\{1,\dots,2017\}$ into the $252$ blocks $\{8k+1,\dots,8k+8\}$ for $k=0,\dots,251$ (covering $1,\dots,2016$) and the single number $2017$. Then
\[
|M|\le2\cdot252+1=505 .
\]
So the maximum is $\boxed{505}$.
